% ---
% title: Personal Website Design
% date: 2026-06
% category:
% nodes: tools
% links:
% ---
\documentclass{article}
\usepackage[utf8]{inputenc}
\usepackage{amsmath, amssymb, amsthm}
\usepackage{geometry}
\usepackage{graphicx}
\usepackage{hyperref}

\geometry{a4paper, margin=1in}

\title{Personal Website Design}
\date{June 2026}
\author{Daksh Mehta}

\begin{document}

Firstly, thank you for visiting my corner of the internet, I hope you stay a while and read a few essays. This essay is for anyone who wants to know about the layout and backend for this website. All code can be found on my \href{https://github.com/daksh-4/personal-website}{github page}. 

This website is custom-built as a collection of static HTML files rather than using a website builder. At its core, it relies on a simple, automated pipeline built with Python that takes raw text and transforms it into the web pages you are reading now. 

\section*{The Writing Environment}
I write my essays as raw text files, typically using \LaTeX (via MiniTex and TexStudio). LaTex was a recently new discovery for me. For those unfamiliar, it is a typesetting system which uses plain code in syntax to give you complete flexibility over your document as opposed to drag-and-drop tools in Word. This is particularly useful for math blocks, integrating figures and custom formatting. I am using it in a relatively simple way but the sky is the limit. 

\section*{The Build Pipeline}
To get these raw text files onto the web, I use a custom Python script that triggers a compiler that translates the LaTex syntax into standard web-readable HTML. Next the script injects this content into a HTML template. The template provides minimal CSS styling and calls MathJax which acts as an engine to handle math symbols. 

\section*{Graph View}
As I write, I tag my essays with core topics and sometimes link them to each other. The web of essay inter-connectedness can be seen on the graph view in the Essays tab. This is a complete rip-off from \href{https://obsidian.md}{Obsidian's} graph view which I have been using for the past 3 years and highly recommend. During the build process, the pipeline compiles essay metadata into a single json, which can be then used by \href{https://github.com/vasturiano/force-graph}{force-graph} to read this data and render an interactive network visualisation. In this view, essays are clustered around their core themes and lines connect related pieces of writing. 

\section*{Hosting and Analytics}
Because the pipeline outputs basic, static HTML, the website requires virtually no computing power to run as there are no databases to query or servers to maintain. I host the site through Porkbun. For tracking viewership, I wanted to avoid invasive cookies. I embedded a tiny tracking code from \href{https://www.goatcounter.com}{GoatCounter}. It is an open-source, free, privacy-friendly analytics tool for smaller websites. It doesn't track users across the web or collect personal information, just logs anonymous pageviews so I can see what people are reading. 

I hope you enjoyed visiting, I found this very satisfying to build. In the future, I might add more features such as mailing lists and a page for downloading some resources I have created so all of this is subject to change. Why not go back to the Home tab to read a random essay? Or check out the graph view in the Essays tab?


\end{document}
